\documentclass[11pt,a4paper]{article}
\usepackage[T1]{fontenc}
\usepackage[utf8]{inputenc}
\usepackage[margin=1in,left=0.75in,right=0.7in,top=0.5in,bottom=0.45in]{geometry}
\usepackage{xcolor}
\usepackage{tikz}
\usepackage{parskip}
\usepackage{enumitem}
\usepackage{needspace}
\usepackage{eso-pic}
\usepackage{graphicx}
\usepackage{listings}
\usepackage{tabularx}
\usepackage{booktabs}
\usepackage[colorlinks=true,urlcolor=teal,linkcolor=teal]{hyperref}

\definecolor{ink}{HTML}{111827}
\definecolor{teal}{HTML}{0F766E}
\definecolor{sand}{HTML}{FBFAF6}
\definecolor{rule}{HTML}{D9D6CC}
\definecolor{muted}{HTML}{6B7280}
\pagecolor{sand}\color{ink}
\pagestyle{empty}
\AddToShipoutPictureBG{\begin{tikzpicture}[remember picture,overlay]
  \fill[teal] (current page.north west) rectangle ([xshift=0.35in]current page.south west);
\end{tikzpicture}}
\newcommand{\kicker}[1]{{\sffamily\bfseries\small\color{teal}\MakeUppercase{#1}}}
\newcommand{\sect}[1]{\needspace{5\baselineskip}\vspace{0.5em}\par\kicker{#1}\par\vspace{-0.5em}{\color{rule}\rule{\linewidth}{0.5pt}}\par\vspace{-0.2em}}
\newcommand{\fig}[2]{\begin{figure}[htbp]\centering\includegraphics[width=\linewidth]{charts/#1}\caption{#2}\end{figure}}
\setlength{\parskip}{0.3em}
\setlist[itemize]{itemsep=1pt,topsep=1pt,leftmargin=1.2em}

\begin{document}
\kicker{ShadowFox data analyst internship}
\begin{flushleft}
{\LARGE\bfseries Azizul Abedin Azmi}\par
{\small\color{muted} \href{mailto:azmi00121@gmail.com}{azmi00121@gmail.com} \quad|\quad \href{https://github.com/azizulabedinazmi}{GitHub} \quad|\quad \href{https://azizulabedinazmi.is-a.dev/}{Portfolio}}
\end{flushleft}
\begin{flushright}
{\sffamily\bfseries\fontsize{25}{28}\selectfont Customer and\\
Revenue Analysis}\\[0.4em]
{\color{muted}\sffamily Online Retail --- Intermediate Level}\\
{\color{muted}\sffamily 11 October 2026}
\end{flushright}

\sect{Executive summary}
This analysis answers where revenue comes from, how well the retailer keeps its
customers, and which actions are most likely to grow revenue. The cleaned
dataset contains 522,537 completed product lines, 19,771 invoices, 4,333
identified customers, 3,908 products and 38 countries. Revenue is
\textbf{\pounds10.00M}; cancellations are \pounds230K, or 2.3\% of revenue.

\begin{tabularx}{\linewidth}{@{}lXr@{}}
\toprule
\textbf{Metric} & \textbf{Definition} & \textbf{Result}\\
\midrule
Revenue & Quantity multiplied by UnitPrice, completed sales & \pounds10.00M\\
Average order value & Revenue divided by invoices & \pounds505.84\\
Repeat customer rate & Customers with at least two orders & 65.3\%\\
Repeat revenue share & Revenue from repeat customers & 93.5\%\\
Top 10\% concentration & Revenue from largest 10\% of identified customers & 60.3\%\\
Identified revenue share & Revenue attached to a CustomerID & 84.9\%\\
\bottomrule
\end{tabularx}

\fig{01_monthly_revenue.png}{Monthly revenue and active customers. December 2011 is partial and ends on 9 December.}

\sect{Business question and method}
The source is the Online Retail dataset from the UCI Machine Learning
Repository, covering 1 December 2010 to 9 December 2011. The grain is one
product line on one invoice. Revenue is \texttt{Quantity x UnitPrice}; no cost
or profit field exists, so every value in this report is revenue.

\begin{itemize}
\item Revenue and product/country analysis uses every valid completed sale,
including lines without a CustomerID.
\item Customer behaviour, retention and RFM use only the 4,333 identified
customers; these represent 84.9\% of revenue.
\item Cohort retention measures the share of each monthly first-order cohort
that orders again after 1, 3 and 6 months. December 2010 is excluded from
averages because the dataset starts then, and December 2011 is partial.
\item RFM scores Recency, Frequency and Monetary value from 1 to 5 by quintile,
then assigns six practical customer segments.
\end{itemize}

\sect{Cleaning and preparation}
The raw export contains 541,909 lines. The cleaning script removes exact
duplicates, separates cancellations, removes non-product charges and invalid
quantity/price lines, standardises product descriptions by StockCode, and
removes two very large sale/cancellation pairs that were reversed in full.
The cleaning log records 15 steps and every row count.

\begin{tabularx}{\linewidth}{@{}lXr@{}}
\toprule
\textbf{Issue} & \textbf{Treatment} & \textbf{Rows}\\
\midrule
Exact duplicate lines & Keep first occurrence & 5,268\\
Cancellations & Move to separate cancellation analysis & 9,251\\
Non-product lines & Remove postage, fees, vouchers and adjustments & 2,356\\
Invalid quantity or price & Remove non-positive normal sales & 2,493\\
Description variants & Use modal description per StockCode & 6,784\\
Missing CustomerID & Keep for revenue; exclude from customer behaviour & 131,388\\
\bottomrule
\end{tabularx}

\sect{Revenue trend, timing and markets}
\fig{02_weekday_hour.png}{Revenue is concentrated in working hours and is strongest from Tuesday to Thursday.}
From January to August, monthly revenue ranges from roughly \pounds0.5M to
\pounds0.75M. It rises to \pounds1.03M in September, \pounds1.10M in October
and \pounds1.45M in November. September to November contributes 37.5\% of
full-month revenue, and November is 2.2 times the January--August average.
The chart shows active customers rising with the peak, not just larger baskets.

About 75.5\% of revenue is placed from 10:00 to 15:59. Thursday and Tuesday
are the strongest days, Sunday contributes 8.0\%, and no Saturday orders are
recorded. This is consistent with trade or business buyers, but the dataset
cannot prove customer type or explain the Saturday gap.

\fig{03_countries.png}{Top export markets by revenue, excluding the United Kingdom.}
The United Kingdom produces 84.8\% of revenue; exports contribute 15.2\%.
The largest export markets are the Netherlands (\pounds284K), EIRE
(\pounds271K), Germany (\pounds205K), France (\pounds184K) and Australia
(\pounds138K). EIRE, the Netherlands and Australia are relatively concentrated
in a small number of customers, while Germany and France are broader markets.

\sect{Customer concentration and repeat purchase}
\fig{04_concentration.png}{Cumulative revenue concentration. The top 10\% of customers produce 60\% of revenue and the top 20\% produce 74\%.}
The top 1\% of customers produce 30.5\% of revenue, the top 10\% produce
60.3\%, and the top 20\% produce 73.9%. The bottom half contributes only 8.1%.
The mean order value is \pounds506 while the median is \pounds302, showing
that a small number of large baskets pull up the average.

\fig{05_orders_per_customer.png}{Distribution of order counts per identified customer.}
Repeat buying is the commercial engine: 65.3\% of identified customers ordered
more than once and repeat customers produce 93.5\% of identified-customer
revenue. Returning-customer baskets average \pounds485 versus \pounds402 for
new-customer order months. The most economical growth lever is therefore
turning first-time buyers into second-time buyers.

\sect{Retention and new versus returning revenue}
\fig{06_cohort_retention.png}{Cohort retention heatmap. Blank cells are unavailable or excluded because December 2011 is partial.}
On average, 19.8\% of a monthly new-customer cohort orders again one month
later, 22.9\% remains active at month three and 25.0\% at month six. Retention
does not decline monotonically because customers reorder in bursts, likely
around their own restocking cycles. The December 2010 cohort is not a true
new-customer cohort and is excluded from the averages.

\fig{07_new_vs_returning.png}{Monthly revenue from new and returning identified customers.}
Returning customers supply 79.7\% of identified-customer revenue in 2011, and
their contribution grows through the year. The November peak is therefore
largely existing-customer expansion rather than only acquisition.

\sect{RFM customer segments}
\fig{08_rfm_segments.png}{RFM segment shares of customers and revenue.}
\begin{tabularx}{\linewidth}{@{}lrrrX@{}}
\toprule
\textbf{Segment} & \textbf{Customers} & \textbf{Customer share} & \textbf{Revenue share} & \textbf{Action}\\
\midrule
Champions & 912 & 21.0\% & 66.2\% & Protect with account ownership and early warning.\\
Loyal & 1,047 & 24.2\% & 15.2\% & Encourage frequency and broader baskets.\\
At Risk & 641 & 14.8\% & 8.6\% & Personal win-back before discounting.\\
Lost / Hibernating & 1,064 & 24.6\% & 5.7\% & Low-cost reactivation tests.\\
Needs Attention & 356 & 8.2\% & 2.3\% & Nurture and observe.\\
New / Promising & 313 & 7.2\% & 2.0\% & Drive the second order.\\
\bottomrule
\end{tabularx}

Champions average 11.4 orders and \pounds6,165 in annual revenue. At Risk
customers average 154 days since their last order and represent \pounds734K of
historic revenue. A narrower high-value quiet cohort contains 246 customers
and \pounds535K of prior revenue, making it a practical first win-back list.

\sect{Products and cancellations}
\fig{09_top_products.png}{Top ten products by revenue. The catalogue is broad: the top ten products account for 8.2\% of revenue, while the top 10\% of products account for 61.0\%.}
The leading product is Regency Cakestand 3 Tier at \pounds174K, followed by
White Hanging Heart T-Light Holder and Party Bunting. The top ten products
represent only 8.2\% of revenue, so availability matters across a wide
catalogue rather than for one dominant item.

Cancellations total \pounds230K, or 2.3\% of revenue. April 2011 is the peak
month at 6.5\%, and the ten customers with the most cancelled value account for
34.2\% of cancellation value. Cancellations are worth monitoring but are not
the first growth priority.

\sect{Recommendations and measurement}
\begin{enumerate}[leftmargin=1.5em,itemsep=2pt]
\item \textbf{Protect Champions.} Assign account owners to the top 10\% and
alert the team when a customer's gap exceeds their normal ordering interval.
Measure Champion retention and order frequency.
\item \textbf{Win back lapsed high-value customers.} Contact At Risk and
valuable quiet customers personally before offering discounts. Test outreach
against a holdout group and measure reactivation revenue.
\item \textbf{Get the second order.} Send a follow-up two to three weeks after
the first purchase with relevant best sellers and a modest incentive. Measure
second-order conversion within 90 days.
\item \textbf{Build Q4 from August.} Forecast products, stock and staffing
before September, then compare Q4 revenue, stock-outs and service levels with
plan.
\item \textbf{Time communications.} Prioritise Tuesday--Thursday mornings and
maintain support coverage from 10:00 to 16:00. Measure conversion by send time.
\item \textbf{Use different export strategies.} Treat EIRE, the Netherlands and
Australia as key-account markets; use broader market marketing in Germany and
France.
\end{enumerate}

\sect{Limitations}
This is one retailer over one year, so seasonality cannot be separated from
growth. Around 25\% of lines have no CustomerID, and customer findings cover
identified customers only. There is no cost, profit, marketing, customer-type
or price data, so explanations are hypotheses rather than causal conclusions.
RFM boundaries are quintile-based and intentionally pragmatic. Cancellations
are analysed as revenue leakage, but the dataset does not provide a reason for
each cancellation.

\sect{Reproducibility}
\begin{lstlisting}
cd "d:\ShadowFox Intership\shadowfox-data-analyst\02-intermediate"
python scripts\clean_online_retail.py
python analysis\analyze_online_retail.py
python analysis\build_report.py
pdflatex -interaction=nonstopmode -halt-on-error -output-directory report report\online_retail_report.tex
\end{lstlisting}
The Excel analysis tables, nine chart PNGs, Word report and this PDF are all
generated from the cleaned CSV outputs.

\end{document}
